\documentclass[../main.tex]{subfiles}

\begin{document}

\section{\textit{Common Analysis File}} \label{sec:common-analysis-file}
  I file di analisi per l'utente finale dell'esperimento DUNE sono mantenuti nel
  formato CAF (\textit{Common Analysis File}).
  I CAF sono file .root che contengono dei \texttt{TTree} basati su degli
  oggetti di tipo \texttt{StandardRecord}.
  Uno \texttt{StandardRecord} è concepito per essere un riassunto
  sufficiente di un evento di neutrino per l'analisi di alto livello.

  In quanto tali, essi riflettono le proprietà di oggetti ricostruiti di alto
  livello (come tracce o sciami), oltre a informazioni sintetiche di verità
  Monte Carlo (\textit{truth}), ma non includono informazioni sulla carica, sui
  singoli \textit{hits} o dettagli temporali di basso livello e sono
  progettati per consentire un'iterazione rapida nei \textit{workflow} di
  analisi; all'interno dei CAF uno \textit{spill} è codificato all'interno di
  un'istanza della classe \texttt{StandardRecord}.
  I CAF sono pensati per avere una bassa barriera all'ingresso e
  per essere facili da usare: gli unici requisiti per leggerli sono ROOT e
  gli header che definiscono gli oggetti di tipo \texttt{StandardRecord}.

\end{document}